\section{The reduction of the charged sector}
\label{app:charged}

\subsection{The zero mode in general dimension}
\label{app:charged-zero-mode}

On the diagonal brane
$ds^2 = -U dt^2 + dr^2/U + e^{2V}(dx^2 + dy^2) + e^{2Z}d\vec z^{\,2}$ in
$d + 1$ dimensions, with $d - 3$ directions $\vec z$ along the field, the
Ricci-form equations are $R_{MN} + d\,g_{MN} = \alpha(F^a_{MP}F^{a\phantom{N}P}_N
- F^2 g_{MN}/2(d-1))$. We write $W = A^1 + iA^2$ and $D = \partial + iA^3$,
so that on $y$-independent functions the Landau ladder is
$D_\pm = D_x \pm iD_y = \partial_x \mp Bx$ with
$D_+\psi_n = -\sqrt B\,\psi_{n+1}$, $D_-\psi_n = 2n\sqrt B\,\psi_{n-1}$ and
$\psi_n = H_n(\sqrt B x)\,e^{-Bx^2/2}$. For the polarised lowest Landau
level $W_+ \equiv W_x + iW_y = p_+(r)\psi_0$, $W_- \equiv W_x - iW_y = 0$, $W_r = 0$, the
linearised Yang--Mills equations reduce, after the $x$ integration is done
exactly, to
\begin{equation}
\big(P w'\big)' + Q w = 0~, \qquad P = U e^{(d-3)Z}~, \qquad Q = B\,e^{(d-3)Z - 2V}~,
\end{equation}
with $w = p_+/2$. The transverse Landau term vanishes because
$F^W_{xy} = \tfrac{i}{2}\big(D_+ W_- - D_- W_+\big)$ has no component in this polarisation, and
$Q$ is the non-abelian moment term $\epsilon^{abc}F^b_{MN}W^{cM}$ with the
single warp factor $e^{-2V}$ of the flux plane. At the same order nothing
sources $A^3$ or the metric, which we verified term by
term, so the truncation is consistent. Near the boundary the two solutions
behave as $w = j_0\big(1 + \tfrac{B}{2}r^{-2}\ln r + \ldots\big) + j_1 r^{-2}
+ \ldots$ in $d = 4$, and the source-free condition is $j_0 = 0$. The
brane's own corrections start at $r^{-3}$ in the fixed-temperature frame,
where $U = r^2 + 2ar + \ldots$, so they are subleading to both terms, and
the source can be read off unambiguously as long as the fit or the
extraction allows for them. At the
horizon, regularity of the solution with $P \sim U'(r_p)(r - r_p)$ requires
$w'(r_p) = -Q(r_p)\,w(r_p)/P'(r_p)$.

\subsection{The other channels}
\label{app:charged-channels}

For the argument of section~\ref{sec:4.5} we keep the
general static, $z$-independent fluctuation in the gauge $W_r = 0$. Since
$D_\pm$ shift the Landau index by one and the background depends on $r$
only, the fluctuations sort by a single integer $n \ge 0$,
\begin{equation}
W_+ = p_+(r)\,\psi_n~, \quad W_- = p_-(r)\,\psi_{n-2}~, \quad W_t = a(r)\,\psi_{n-1}~, \quad W_z = c(r)\,\psi_{n-1}~,
\end{equation}
with $\psi_{k<0} \equiv 0$. Integrating the quadratic action over $x$ is
exact at every $n$ and produces no coupling between different $n$. The
components $W_t$ and $W_z$ decouple from $(p_+, p_-)$ and from each other, each
with a Sturm--Liouville operator $(P_c a')' - (2n-1)B\zeta_c\,a$ whose
coefficients $P_c, \zeta_c$ are positive and whose transverse mass
$(2n-1)B$ is positive for $n \ge 1$. In the convention of appendix~\ref{app:charged-zero-mode}, with
$P = \sqrt{-g_{tt}g_{zz}/g_{rr}}$, $\zeta = \sqrt{-g_{tt}g_{rr}g_{zz}}/g_{xx}$
(on the brane $P = p_2$ and $\zeta = p_1 = Q/B$, the weights of
section~\ref{sec:5}) and $N_k = \int\psi_k^2\,dx$, the remaining density, written as minus the
energy so that a positive $T_n$ is destabilising, is
\begin{equation}
\mathcal{E}_n = -\frac{P}{4}\big(N_n\,p_+'^2 + N_{n-2}\,p_-'^2\big) + T_n(p_+, p_-)~, \qquad
T_0 = +\frac{B\zeta N_0}{4}\,p_+^2~, \qquad T_1 = 0~,
\end{equation}
and, for $n \ge 2$,
\begin{equation}
\label{eq:Tn}
T_{n} = -\frac{B\zeta N_{n-1}}{4}\begin{pmatrix} 2n(n-1) & n \\ n & \dfrac{n}{2(n-1)} \end{pmatrix}~,
\end{equation}
which has vanishing determinant and negative trace. There are no cross
terms between derivatives and undifferentiated fields, and the real and
imaginary parts of $(p_+, p_-)$ carry identical blocks. The null vector of
$T_{n\ge2}$, $(p_+, p_-) \propto (-1, 2(n-1))$, is the residual gauge
transformation $W_M = D_M(\theta\psi_{n-1})$, as the ladder relations
$D_+\psi_{n-1} = -\sqrt B\psi_n$ and $D_-\psi_{n-1} = 2(n-1)\sqrt B\psi_{n-2}$
show, and with $W_r$ reinstated the density shifts by a total derivative
under it, which we used as the check on the reduction. For general $n$ the transverse form follows from the ladder relations
alone: the transverse energy is $\tfrac12\zeta$ times
$\tfrac14|D_+W_- - D_-W_+|^2 - \tfrac B2\big(|W_+|^2 - |W_-|^2\big)$, and with $D_+\psi_{n-2} = -\sqrt B\psi_{n-1}$,
$D_-\psi_n = 2n\sqrt B\psi_{n-1}$ and $N_n = 2nN_{n-1}$ the $x$ integral
gives $\tfrac{B\zeta N_{n-1}}{4}\big[2n(n-1)p_+^2 + 2np_+p_- + \tfrac{n}{2(n-1)}p_-^2\big]$,
which is $-T_n$ of~\eqref{eq:Tn}, as a symbolic reduction for
$n = 0, \ldots, 4$ on the general diagonal metric confirms.
The statement of section~\ref{sec:4.5} follows: for $n \ge 1$ the energy is a sum of
positive radial terms and a positive semi-definite transverse form, and it
vanishes only on $r$-independent pure gauge, so a source-free, horizon-regular
zero mode must be pure gauge, while $n = 0$, with $p_+ = 2w$, reproduces
$N_0$ times the density $-Pw'^2 + Qw^2$, whose Euler--Lagrange equation
is the zero-mode equation of appendix~\ref{app:charged-zero-mode}.

The argument of section~\ref{sec:4.5} is made for static, $z$-independent
perturbations. With a dependence $e^{ik_zz}$ the $(x,z)$ and $(y,z)$ terms
of the energy become
$\tfrac14\sqrt{-g}\,g^{xx}g^{zz}\big(|D_+W_z - ik_zW_+|^2 + |D_-W_z - ik_zW_-|^2\big)$,
two squares that mix $W_z$ with $W_\pm$, while the $(x,y)$ form is unchanged.
For $n \ge 1$ the energy therefore stays non-negative, with kernel the pure
gauge $D_M(\theta\psi_{n-1}e^{ik_zz})$, and for $n = 0$, where
$W_- = W_z = 0$, the new term is the positive mass $N_0 e^{-Z}k_z^2w^2$.
A dependence on $z$ cannot produce a zero mode where the $k_z = 0$ operator
has none.

The onset is also static. The background has no $A_t$ and no $g_{ti}$, so
time derivatives enter the quadratic action only through $|F_{ti}|^2$:
\begin{equation}
L = \tfrac12\sqrt{-g}\,(-g^{tt})g^{ii}|F_{ti}|^2 - \mathcal V[W]~,
\end{equation}
where $\mathcal V$ is the static energy above. The magnetic field makes the
operator of $\mathcal V$ Hermitian rather than real symmetric, but it does
not add a term linear in $\partial_t$. Take a mode $e^{-i\omega t}$ with
$\mathrm{Im}\,\omega > 0$, ingoing and source-free, in the gauge $W_t = 0$.
The components of $W$ in this gauge then vanish at the horizon like
$(r - r_p)^{\mathrm{Im}\,\omega/U'(r_p)}$; $F_{ri}$ grows like
$(r - r_p)^{\mathrm{Im}\,\omega/U'(r_p) - 1}$, but every integrand goes
like $(r - r_p)^{2\mathrm{Im}\,\omega/U'(r_p) - 1}$ and every surface term
like $(r - r_p)^{2\mathrm{Im}\,\omega/U'(r_p)}$. Contracting the equations
with $\bar W$ therefore gives, without surface terms,
\begin{equation}
\frac{\omega^2}{|\omega|^2}\int\sqrt{-g}\,(-g^{tt})g^{ii}|F_{ti}|^2 = 2\,\mathcal V[W]~.
\end{equation}
Hence $\omega^2$ is real, so $\omega = i\Gamma$, and the static energy of
the mode is negative. For $n \ge 1$, $\mathcal V \ge 0$ and no mode grows,
also at $k_z \neq 0$: $W_r$, which this gauge leaves in place, enters
$\mathcal V$ only through the squares $|F_{ri}|^2$, and the transverse form
does not contain it. For $n = 0$ the mode obeys
$(Pw')' + Qw = \Gamma^2 R\,w$ with $R = e^{Z}/U$ positive,
a self-adjoint problem whose continuum begins at $\Gamma^2 = 0$; by Sturm
oscillation the number of its growing modes equals $\#\{k : B_k < B\}$,
and they emerge through $\omega = 0$. A chemical potential would add a term linear in
$\partial_t$ and shift $\omega$ to $\omega - \mu$, which is where the
argument stops.

We checked the $n = 0$, $k_z = 0$ channel numerically on the brane at
$\beta = 1$, $45$ and $10^4$ and in the probe limit, by collocation in
ingoing coordinates and independently by the argument principle applied to
the shooting source coefficient. Both count exactly
$\#\{k : B_k < B\}$ modes in the upper half plane; the collocation
eigenvalues are purely imaginary to
$2\times10^{-11}$, and one of them crosses $\omega = 0$ exactly at $B_c$.
